% Source: https://artofproblemsolving.com/community/c3h2409265_a_quick_question_on_area
% Stable identity: aops-topic-2409265-q01
\textbf{The locus of $C$.} Put $A=(0,0)$, $B=(9,0)$ and $C=(x,y)$ with $y>0$. The condition $41\cdot BC=40\cdot AC$ means
\[1681\left((x-9)^2+y^2\right)=1600\left(x^2+y^2\right)\iff81x^2+81y^2-30258x+136161=0.\]
Dividing by $81$ and completing the square,
\[x^2-\dfrac{3362}{9}x+y^2+1681=0\iff\left(x-\dfrac{1681}{9}\right)^2+y^2=\left(\dfrac{1681}{9}\right)^2-1681=\dfrac{1681\cdot1600}{81}.\]
So $C$ lies on the circle with centre $O=\left(\dfrac{1681}{9},0\right)$ and radius $\dfrac{41\cdot40}{9}=\dfrac{1640}{9}$ (a circle of Apollonius).

\textbf{The maximum.} The area is $\dfrac12\cdot AB\cdot y=\dfrac92y$, and on this circle $y\le\dfrac{1640}{9}$, with equality at the top point $C=\left(\dfrac{1681}{9},\dfrac{1640}{9}\right)$, which gives a genuine triangle with $BC:AC=40:41$. Hence the largest area is
\[\dfrac92\cdot\dfrac{1640}{9}=\boxed{820},\]
and the sum of its digits is $8+2+0=\boxed{10}$.

\begin{asy}
usepackage("amsmath");
size(6cm);
defaultpen(fontsize(10pt));
real s=1/40;
pair A=(0,0), B=(9,0), O=(1681/9,0), C=(1681/9,1640/9), H=(1681/9,0);
draw(scale(s)*circle(O,1640/9),dashed+gray(0.5));
draw(scale(s)*(A--B--C--cycle),blue+linewidth(0.9));
draw(scale(s)*(C--H),dashed);
draw(scale(s)*((-12,0)--(380,0)),gray(0.6));
dot(scale(s)*A); dot(scale(s)*B); dot(scale(s)*C); dot(scale(s)*O);
label("$A$",scale(s)*A,SW); label("$B$",scale(s)*B,SE); label("$C$",scale(s)*C,N); label("$O$",scale(s)*O,S);
label("$\dfrac{1640}{9}$",scale(s)*(C+H)/2,E);
\end{asy}

The figure is to scale: $AB=9$ is tiny compared with the circle, and the best triangle is long and thin ($AC=41k$, $BC=40k$ with $k=\dfrac{\sqrt{3281}}{9}\approx6.36$).
